\chapter{No convergencia}

\section{Un gran debate}

En el siglo XIX, cuando la gente aún tenía confusión sobre los límites en general,
existía una sucesión concreta sobre la que se discutía acaloradamente:
\[
  1,\; -1,\; 1,\; -1,\; 1,\; -1,\; \ldots
\]
¿Converge o no converge?

Algunos decían: subes uno y bajas uno, así que en promedio está alrededor de
cero. Por tanto, el límite debería ser~$0$.

Otros respondían: fíjate en las sumas parciales de la sucesión:
\begin{align*}
  1 &= 1 \\
  1 + (-1) &= 0 \\
  1 + (-1) + 1 &= 1 \\
  1 + (-1) + 1 + (-1) &= 0
\end{align*}
Las sumas parciales alternan entre~$1$ y~$0$, así que quizá el límite sea
$\tfrac{1}{2}$.

Esto es el tipo de problema al que te enfrentas cuando no existe un consenso
comunitario sobre las definiciones, es decir, cuando no hay una definición
rigurosa de límite. ¡Por suerte, ahora sí la tenemos!

Con esa definición, esta sucesión \ldots\ \textbf{no tiene límite} (diverge).

Para expresarlo formalmente, introduzcamos una nueva definición:

\begin{minted}{lean4}
def SucConv (a : ℕ → ℝ) : Prop :=
  ∃ L, LimSuc a L
\end{minted}

Es decir, una sucesión~$a$ «converge» (sin especificar a qué) si existe algún
número real~$L$ tal que $a$ converge a ese número.

El objetivo de este capítulo es demostrar:
\[
  \text{Meta}: \neg\,\texttt{SucConv}\;a,
\]
dado la hipótesis
\[
  \texttt{ha : ∀n, a n = (-1 : R)\^{}n}
\]
(Nota: el símbolo $\neg$ se obtiene escribiendo \verb|\not|.)

¿Cómo se demuestra la negación de algo? Lógicamente hablarías así: si eso
ocurriera, obtendríamos una contradicción. Técnicamente, $\neg P$ es
por definición equivalente a \[P \to \text{False}\]

Por tanto, podemos comenzar la demostración recordándonos este hecho escribiendo:
\begin{minted}{lean4}
change SucConv a → False
\end{minted}
Eso cambiará la meta a: \texttt{SucConv a → False}.

En general, recomiendo el siguiente protocolo:

\begin{enumerate}
  \item \textbf{Pensar.} Hacer un borrador, tener una idea vaga de cómo
        ensamblar el argumento.
  \item \textbf{Demostrarlo formalmente.} ¡Pero no has terminado! Que algo
        funcione formalmente no significa necesariamente que entiendas qué está
        ocurriendo. El objetivo de las matemáticas \emph{no} es que simplemente
        funcione, sino que \textbf{entiendas exactamente} qué está pasando y por
        qué. Por eso existe el paso~3:
  \item \textbf{Dar una demostración en lenguaje natural} que explique
        (para ti mismo tanto como para cualquier otra persona, incluyéndome a
        mí) qué ocurre y por qué.
\end{enumerate}

Usando únicamente reacciones instintivas (si ves $P \to Q$ en la meta,
escribe \texttt{intro}; si ves $\exists$ en una hipótesis, escribe
\texttt{choose}; si ves una definición en la meta o en una hipótesis,
escribe \texttt{change}), llegamos al siguiente punto:

\medskip
\noindent\textbf{Objetos:} $a : \mathbb{N} \to \mathbb{R}$, $L : \mathbb{R}$

\noindent\textbf{Supuestos:}
\begin{itemize}
  \item \texttt{ha : ∀(n : N), a n = (-1)\^{}n}
  \item \texttt{hL : ∀ε > 0, ∃N, ∀n ≥ N, |a n - L| < ε}
\end{itemize}

\noindent\textbf{Meta:} \texttt{False}

\medskip
Recuerda al Ingeniero y al Mecánico: el Ingeniero especifica cualquier
tolerancia, y el Mecánico debe garantizar que estaremos dentro de esa
tolerancia «eventualmente», es decir, para todo~$n$ suficientemente grande,
medido por la cota inferior~$N$.

\medskip
\noindent\textbf{Idea clave:} Encontrar un~$\varepsilon$ cuya tolerancia sea
demasiado ajustada para que el Mecánico pueda cumplirla. Como la sucesión
alterna entre~$1$ y~$-1$, que difieren en~$2$, la tolerancia $\varepsilon = 1$
ya sería suficiente, pero para tener un poco más de margen tomemos
$\varepsilon = \tfrac{1}{2}$.

\begin{nota}
  Si hubiéramos tomado $\varepsilon = 1.5$, \emph{no} funcionaría. No podemos
  descartar la posibilidad de que $L = 0$, y los valores~$1$ y~$-1$ \emph{sí}
  están dentro de esa tolerancia. Por tanto, no hallaríamos ninguna
  contradicción.
\end{nota}

Tras varios cálculos, llegamos a:
\[
  |1 - L| < \tfrac{1}{2}
  \qquad\text{y}\qquad
  |-1 - L| < \tfrac{1}{2}.
\]
Intentemos obtener una contradicción a partir de ahí:
\begin{align*}
  2 &= |2| \\
    &= |1 - (-1)| \\
    &= |(1 - L) + (L - (-1))| \\
    &\leq |(1 - L)| + |(L - (-1))| \\
    &= |(1 - L)| + |{-}((-1) - L)| \\
    &= |(1 - L)| + |((-1) - L)| \\
    & < \tfrac{1}{2} + \tfrac{1}{2} \\
    &= 1.
\end{align*}
Eso nos da la contradicción deseada.

% ─── Nuevas herramientas ──────────────────────────────────────────────────────
\subsection*{Nuevas herramientas que necesitarás}

\begin{itemize}
  \item \textbf{Negación:} Si $P$ es alguna \texttt{Prop}, entonces $\neg P$
        es por definición equivalente a $P \to \text{False}$. Por tanto,
        puedes escribir `\texttt{change P → False}`, ya sea en la meta o en una
        hipótesis.

  \item \textbf{La desigualdad triangular:} Para añadir el hecho de que
        $|x + y| \leq |x| + |y|$ a nuestra lista de hipótesis, invoca el
        teorema \texttt{abs\_add}:
\begin{minted}{lean4}
have h : |x + y| ≤ |x| + |y| := by apply abs_add
\end{minted}

  \item \textbf{Negación dentro de un valor absoluto:} También puede
        resultar útil el teorema \texttt{abs\_neg}, que se llama así:
\begin{minted}{lean4}
have h : |-x| = |x| := by apply abs_neg
\end{minted}

  \textbf{¡Atención!} Asegúrate de que el patrón \texttt{|-Algo|} esté en el
  lado izquierdo. Si Lean no ve un valor absoluto seguido de un signo menos,
  \texttt{abs\_neg} no funcionará.
\end{itemize}

% ══════════════════════════════════════════════════════════════════════════════
\section{Paso 2: Solución en Lean}

\begin{minted}{lean4}
-- Demostración en lenguaje natural
-- ================================

-- Sea a la sucesión definida por a(n) = (-1)^n. Supongamos que a es
-- convergente. Entonces, existe un L tal que a converge L. Por tanto,
-- existe un N ∈ ℕ tal que,
--    ∀n ≥ N, |a(n) - L| < 1/2                                      (1)
-- Entonces,
--    2 = |2|
--      = |(1 - L) + (1 + L)|
--      = |(1 - L) + (-1)(-1 - L)|
--      ≤ |1 - L| + |(-1)(-1 - L)|
--      = |1 - L| + |-1 - L|
--      = |(-1)^{2N} - L| + |(-1)^{2N+1} - L|
--      = |a(2N) - L| + |a(2N+1) - L|
--      < 1/2 + 1/2                            [por (1), 2N ≥ N y 2N+1 ≥ N]
--      = 1
-- Luego, 2 < 1 que es una contradicción.

-- 1ª demostración
-- ===============

example
  (ha : ∀ n, a n = (-1) ^ n)
  : ¬ SucConv a :=
by
  intro h
  -- h : SucConv a
  -- ⊢ False
  choose L hL using h
  -- L : ℝ
  -- hL : LimSuc a L
  choose N hN using hL (1/2) (by grind)
  -- N : ℕ
  -- hN : ∀ n ≥ N, |a n - L| < 1 / 2
  have h1 : (2:ℝ) < 1 := by calc
    2 = |2|                                        := by grind
    _ = |(1 - L) + (1 + L)|                        := by grind
    _ = |(1 - L) + (-1)*(-1 - L)|                  := by grind
    _ ≤ |1 - L| + |(-1)*(-1 - L)|                  := by grind
    _ = |1 - L| + |-1 - L|                         := by grind
    _ = |(-1:ℝ)^(2*N) - L| + |(-1:ℝ)^(2*N+1) - L|  := by
          have h2 : (-1:ℝ)^(2*N) = 1    := by simp
          have h3 : (-1:ℝ)^(2*N+1) = -1 := by grind
          rw [h2, h3]
    _ = |a (2*N) - L| + |a (2*N+1) - L|            := by simp [ha]
    _ < 1/2 + 1/2                                  := by grind
    _ = 1                                          := by grind
  linarith

-- 2ª demostración
-- ===============

variable {x y z x' y' : ℝ}
variable {m n k : ℕ}
variable (f : ℝ → ℝ)

lemma L1 : (1 - x) + (1 + x) = 2 := by
  calc (1 - x) + (1 + x)
       = 1 + 1             := sub_add_add_cancel 1 1 x
     _ = 2                 := one_add_one_eq_two

lemma L2 : (-1:ℝ) * -1 = 1 := by
  calc (-1:ℝ) * -1
       = 1 * 1   := neg_mul_neg 1 1
     _ = 1       := one_mul 1

lemma L3 : -1 * -x = x := by
  calc -1 * -x
       = -(1 * -x) := neg_mul 1 (-x)
     _ = - -x      := congr_arg (- ·) (one_mul (-x))
     _ = x         := neg_neg x

lemma L4 : -1 * (-1 - x) = 1 + x := by
  calc -1 * (-1 - x)
       = -1 * (-1 + -x)     := congr_arg (-1 * ·) (sub_eq_add_neg (-1) x)
     _ = -1 * -1 + -1 * -x  := left_distrib (-1) (-1) (-x)
     _ = 1 + -1 * -x        := congr_arg (· + (-1)*(-x)) L2
     _ = 1 + x              := congr_arg ( 1 + ·) L3

lemma L5 : |-1 * (-1 - x)| = |-1 - x| := by
  calc |-1 * (-1 - x)|
       = |-1| * |-1 - x| := abs_mul (-1) (-1 - x)
     _ = |1| * |-1 - x|  := congrArg (· * |-1 - x|) (abs_neg 1)
     _ = 1 * |-1 - x|    := congrArg (· * |-1 - x|) abs_one
     _ = |-1 - x|        := one_mul |-1 - x|

lemma L6 : (-1:ℝ)^(2*n) = 1 := by
  calc (-1:ℝ)^(2*n)
       = ((-1:ℝ)^2)^n := pow_mul (-1) 2 n
     _ = (1:ℝ)^n      := congr_arg ( · ^ n) neg_one_sq
     _ = 1            := one_pow n

lemma L7 : (-1:ℝ)^(2*n+1) = -1 := by
  calc (-1:ℝ)^(2*n+1)
       = (-1)^(2*n) * -1 := pow_succ (-1) (2 * n)
     _ = 1 * -1          := congr_arg (· * -1) L6
     _ = -1              := one_mul (-1)

lemma L8 : n ≤ 2 * n := by
  calc n
       = 1 * n := (one_mul n).symm
     _ ≤ 2 * n := Nat.mul_le_mul_right n one_le_two

lemma L9
  (hN : ∀ n ≥ N, |a n - x| < 1 / 2)
  : |a (2*N) - x| < 1 / 2 :=
hN (2*N) L8

lemma L10 : n ≤ 2 * n + 1 := by
  calc n
       ≤ 2 * n     := L8
     _ ≤ 2 * n + 1 := Nat.le_add_right (2 * n) 1

lemma L11
  (hN : ∀ n ≥ N, |a n - x| < 1 / 2)
  : |a (2*N+1) - x| < 1 / 2 :=
hN (2*N+1) L10

example
  (ha : ∀ n, a n = (-1) ^ n)
  : ¬ SucConv a :=
by
  intro h
  -- h : SucConv a
  -- ⊢ False
  choose L hL using h
  -- L : ℝ
  -- hL : LimSuc a L
  choose N hN using hL (1/2) one_half_pos
  -- N : ℕ
  -- hN : ∀ n ≥ N, |a n - L| < 1 / 2
  have h1 : (2:ℝ) < 1 := by calc
    2 = |2| :=
          abs_two.symm
    _ = |(1 - L) + (1 + L)| :=
          congrArg abs L1.symm
    _ = |(1 - L) + (-1)*(-1 - L)| :=
          congr_arg (abs ((1 - L) + ·)) L4.symm
    _ ≤ |1 - L| + |(-1)*(-1 - L)| :=
          abs_add_le (1 - L) (-1 * (-1 - L))
    _ = |1 - L| + |-1 - L| :=
          congr_arg (|1 - L| + ·) L5
    _ = |(-1:ℝ)^(2*N) - L| + |(-1:ℝ)^(2*N+1) - L|  :=
          congrArg₂ (|· - L| + |· - L|) L6.symm L7.symm
    _ = |a (2*N) - L| + |a (2*N+1) - L| :=
          congrArg₂ (· + ·)
            (congrArg (|· - L|) (ha (2*N)).symm)
            (congrArg (|· - L|) (ha (2*N+1)).symm)
    _ < 1 / 2 + 1 / 2 :=
          add_lt_add (L9 hN) (L11 hN)
    _ = 1 := add_halves 1
  have h2 : (1:ℝ) < 1 := lt_trans one_lt_two h1
  exact (lt_irrefl 1) h2

-- Lemas usados
-- ============

#check (Nat.le_add_right n m : n ≤ n + m)
#check (Nat.mul_le_mul_right k : n ≤ m → n * k ≤ m * k)
#check (abs_add_le x y : |x + y| ≤ |x| + |y|)
#check (add_halves x : x / 2 + x / 2 = x)
#check (add_lt_add : x < x' → y < y' → x + y < x' + y')
#check (congrArg₂  (· + ·) : x = x' → y = y' → x + y = x' + y')
#check (congr_arg f  : x = y → f x = f y)
#check (left_distrib x y z : x * (y + z) = x * y + x * z)
#check (lt_irrefl x : ¬x < x)
#check (lt_trans : x < y → y < z → x < z)
#check (neg_mul 1 x : -1 * x = -(1 * x))
#check (neg_mul_neg x y : -x * -y = x * y)
#check (neg_neg x : - -x = x)
#check (neg_one_sq : (-1)^2 = 1)
#check (one_add_one_eq_two : 1 + 1 = 2)
#check (one_le_two : 1 ≤ 2)
#check (one_lt_two : (1 : ℝ) < 2)
#check (one_mul x : 1 * x = x)
#check (one_mul x : 1 * x = x)
#check (one_pow n : 1 ^ n = 1)
#check (pow_mul x m n : x ^ (m * n) = (x ^ m) ^ n)
#check (pow_succ x n : x ^ (n + 1) = x ^ n * x)
#check (sub_add_add_cancel x y z : (x - z) + (y + z) = x + y)
#check (sub_eq_add_neg x y : x - y = x + -y)

-- 3ª demostración
-- ===============

example
  (ha : ∀ n, a n = (-1) ^ n)
  : ¬ SucConv a :=
by
  change SucConv a → False
  -- ⊢ SucConv a → False
  intro h
  -- h : SucConv a
  -- ⊢ False
  change ∃ L, LimSuc a L at h
  -- h : ∃ L, LimSuc a L
  choose L hL using h
  -- L : ℝ
  -- hL : LimSuc a L
  change ∀ ε > 0, ∃ N, ∀ n ≥ N, |a n - L| < ε at hL
  -- hL : ∀ ε > 0, ∃ N, ∀ n ≥ N, |a n - L| < ε
  specialize hL (1/2)
  -- hL : 1 / 2 > 0 → ∃ N, ∀ n ≥ N, |a n - L| < 1 / 2
  have h0 : (0 : ℝ) < 1/2 := by norm_num
  specialize hL h0
  -- hL : ∃ N, ∀ n ≥ N, |a n - L| < 1 / 2
  choose N hN using hL
  -- N : ℕ
  -- hN : ∀ n ≥ N, |a n - L| < 1 / 2
  have h2N : N ≤ 2 * N := by bound
  have h2Np1 : N ≤ 2 * N + 1 := by bound
  have h1 : |a (2 * N) - L| < 1 / 2 := by apply hN (2 * N) h2N
  have h2 : |a (2 * N + 1) - L| < 1 / 2 := by apply hN (2 * N + 1) h2Np1
  have h3 : a (2 * N) = (-1) ^ (2 * N) := by apply ha (2 * N)
  have h4 : a (2 * N + 1) = (-1) ^ (2 * N + 1) := by apply ha (2 * N + 1)
  rewrite [h3] at h1
  -- h1 : |(-1) ^ (2 * N) - L| < 1 / 2
  rewrite [h4] at h2
  -- h2 : |(-1) ^ (2 * N + 1) - L| < 1 / 2
  have h5 : (-1 : ℝ) ^ (2 * N) = 1 := by bound
  rewrite [h5] at h1
  -- h1 : |1 - L| < 1 / 2
  have h6 : (-1 : ℝ) ^ (2 * N + 1) = -1 := by grind
  rewrite [h6] at h2
  -- h2 : |-1 - L| < 1 / 2
  have h7 : (2 : ℝ) = |2| := by norm_num
  have h8 : |(2 : ℝ)| = |1 - (-1)| := by ring_nf
  have h9 : |1 - (-1)| = |(1 - L) + (L - (-1))| := by noncomm_ring
  have h10 : |(1 - L) + (L - (-1))| ≤ |(1 - L)| + |(L - (-1))| :=
    abs_add_le (1 - L) (L - -1)
  have h11 : |(L - (-1))| = |-((-1) - L)| := by ring_nf
  have h12 : |-((-1) - L)| = |((-1) - L)| := by apply abs_neg
  have h13 : (2 : ℝ) < 1/2 + 1/2 := by
    linarith [h8, h9, h10, h11, h12, h7, h1, h2]
  norm_num at h13
\end{minted}

% ══════════════════════════════════════════════════════════════════════════════
\section{Paso 3: Demostración en lenguaje natural}

\subsection{Demostración en lenguaje natural}

\begin{theorem}
  La sucesión $1, -1, 1, -1, \ldots$ no converge.
\end{theorem}

\begin{proof}[Demostración por contradicción]
\mbox{}
\begin{enumerate}
  \item Supongamos que la sucesión converge a algún límite~$L$.
  \item Tomamos $\varepsilon = \tfrac{1}{2}$ y aplicamos la definición de
        convergencia.
  \item Esto garantiza la existencia de algún~$N$ tal que
        $|a_n - L| < \tfrac{1}{2}$ para todo $n \geq N$.
  \item Consideramos dos índices concretos: $n = 2N$ y $n = 2N+1$
        (ambos son $\geq N$):
        \begin{itemize}
          \item $a(2N) = (-1)^{2N} = 1$, luego $|1 - L| < \tfrac{1}{2}$.
          \item $a(2N+1) = (-1)^{2N+1} = -1$, luego $|-1 - L| < \tfrac{1}{2}$.
        \end{itemize}
  \item Pero entonces:
        \begin{align*}
          2 &= |1 - (-1)| \\
            &= |(1 - L) + (L + 1)| \\
            &\leq |1 - L| + |L + 1| \quad \text{(desigualdad triangular)} \\
            &= |1 - L| + |-1 - L| \quad \text{(manipulación algebraica)} \\
            &< \tfrac{1}{2} + \tfrac{1}{2} \\
            &= 1 \quad \text{(de los pasos anteriores).}
        \end{align*}
  \item Esto nos da $2 < 1$, lo cual es imposible. \qedhere
\end{enumerate}
\end{proof}

\subsection{Ideas clave}

\begin{itemize}
  \item \textbf{Intuición geométrica:} Cualquier límite propuesto~$L$ debe
        estar simultáneamente a distancia menor de $\tfrac{1}{2}$ tanto de~$1$
        como de~$-1$, pero estos valores están a distancia~$2$ entre sí.

  \item \textbf{Elección crítica:} Escogimos $\varepsilon = \tfrac{1}{2}$
        de forma estratégica; valores mayores como $\varepsilon = 1.5$ no
        funcionarían, ya que $L = 0$ podría satisfacer las restricciones.

  \item \textbf{Desigualdad triangular:} El paso clave usa
        $|a + b| \leq |a| + |b|$ para convertir la distancia entre los valores
        de la sucesión en una suma de distancias al límite.
\end{itemize}

Esta demostración ilustra cómo las definiciones rigurosas resuelven debates
matemáticos históricos y proporcionan criterios claros para la convergencia.

%%% Local Variables:
%%% mode: LaTeX
%%% TeX-master: "Analisis_en_Lean"
%%% End:
